% Included in the historical opening subsection of Section 9.
\paragraph{Radial normal forms and the classical model families.}
\label{rem:cubic-radial-normal-form-history}
Choosing a positive maximum direction on the unit sphere, with
coordinate $z$, and normalizing $F$ gives
\[
 \begin{gathered}
 F=z^3+zB+C,\qquad
 B=-3|\xi|^2-\tfrac32|\eta|^2+\tfrac32|\zeta|^2,\\
 C=C_{111}+C_{102}+C_{012}+\sigma,\qquad \sigma=C_{030}.
 \end{gathered}
\]
The subscripts denote multidegrees in $(\xi,\eta,\zeta)$, whose
dimensions are $n_1,n_2,n_3$, respectively. They satisfy
$n_3=2n_1+n_2-2$ and
$N=3n_1+2n_2-1$; see \cite[Section~1, Theorem~1]{PengXiao93}
and \cite[Proposition~3.1]{TkachevClassification10}.

The component $\sigma$ satisfies
$|\nabla\sigma|^2=\tfrac92|\eta|^4$.
The zero case means that the $\eta$-space is absent ($n_2=0$).
Otherwise Peng--Xiao \cite[Theorem~2]{PengXiao93} classify $\sigma$
as reducible or a scaled Cartan cubic; a different proof of the
eikonal classification is given in
\cite[Theorem~1.1]{TkachevCartan10}.
The case $\sigma=0$ yields the Cartan cubics, whereas the case
$\sigma\ne0$ with $\sigma$ reducible yields the Clifford cubics
\cite[Theorems~A and~B]{PengXiao93}; see also
\cite[Proposition~3.4 and Theorem~3]{TkachevClassification10}.
In the remaining branch, $n_2\in\{5,8,14,26\}$.
The cases $n_1=0$ and $n_1=1$ were treated in
\cite[Theorems~C and~D]{PengXiao93} and
\cite[Sections~6.2--6.3]{TkachevClassification10}:
the former yields the three models invariant under $SO(4)$, $SU(4)$,
and $Sp(4)$, and the latter yields four models interpreted through
complexified rank-three Jordan norms
\cite[Section~6.11.6(iii)--(iv)]{NadirashviliTkachevVladut14}.

Tkachev also identifies a hidden Clifford system encoded by
$C_{111}+3^{-1/2}C_{102}$ when $n_1>0$, imposing Hurwitz--Radon restrictions on the
multiplicities \cite[Proposition~3.3]{TkachevClassification10}.
His construction of a $G_2$-invariant cubic of type $(4,5)$ and his
correspondence with isoparametric quartics provide further examples and
restrictions \cite[Section~6.4 and Theorem~5]{TkachevClassification10},
using the isoparametric classification results of
Cecil--Chi--Jensen \cite{CecilChiJensen07} and Immervoll
\cite{Immervoll08}. The exceptional branch has only finitely many
possible multiplicity types, with $N\le72$
\cite[Theorem~6.11.2]{NadirashviliTkachevVladut14}.
This bounds the possible types rather than completing the
classification of their realizations.

